\documentclass[border=10pt]{standalone}
\usepackage{fontspec}
\setmainfont{Noto Serif}[Ligatures=TeXOff]
\XeTeXgenerateactualtext=1
\newfontfamily\bengalifont{Noto Serif Bengali}[Script=Bengali, NFSSFamily=bnfam]
\usepackage[Latin,Bengali]{ucharclasses}
% group-free font switching (as in header-common.tex): robust inside braces, cells and quotes
\makeatletter
\def\BN@saved{\rmdefault}
\def\BN@bn{bnfam}
\newcommand\BNin{\edef\BN@tmp{\f@family}\ifx\BN@tmp\BN@bn\else\let\BN@saved\BN@tmp\fi\fontfamily{bnfam}\selectfont}
\newcommand\BNout{\fontfamily{\BN@saved}\selectfont}
\makeatother
\setTransitionsFor{Bengali}{\BNin}{\BNout}
% Typographic quotes and dashes are not in a ucharclasses class (class 0), so a Bangla word right
% after “ or ‘ would not trigger the Bangla font. Treat them as Latin so that every switch fires.
\XeTeXcharclass"2018=\XeTeXcharclass`A
\XeTeXcharclass"2019=\XeTeXcharclass`A
\XeTeXcharclass"201C=\XeTeXcharclass`A
\XeTeXcharclass"201D=\XeTeXcharclass`A
\XeTeXcharclass"2013=\XeTeXcharclass`A
\XeTeXcharclass"2014=\XeTeXcharclass`A
\XeTeXcharclass"2026=\XeTeXcharclass`A
\usepackage{newunicodechar}
\newunicodechar{।}{{\bengalifont\char"0964}}
\newunicodechar{॥}{{\bengalifont\char"0965}}

\usepackage[table]{xcolor}
\usepackage{booktabs,tabularx,array}
\definecolor{bapilrule}{HTML}{1F2937}
\definecolor{bapilstripe}{HTML}{F3F4F6}
\renewcommand{\arraystretch}{1.28}
\setlength{\tabcolsep}{5pt}

\begin{document}
\begin{minipage}{17cm}
{\large\bfseries সারণি ৪২।}\ {\large BAPIL-এর সঙ্গে প্রাসঙ্গিক বাংলাদেশের আইন ও নীতি-দলিল (প্রতিবেদিত অবস্থা, অক্টোবর ২০২৬)}\par\vspace{7pt}
\small
\rowcolors{2}{bapilstripe}{white}
\begin{tabularx}{\linewidth}{@{}>{\hsize=0.735\hsize\raggedright\arraybackslash}X>{\hsize=1.347\hsize\raggedright\arraybackslash}X>{\hsize=1.020\hsize\raggedright\arraybackslash}X>{\hsize=0.898\hsize\raggedright\arraybackslash}X@{}}
\arrayrulecolor{bapilrule}\toprule
\rowcolor{white}\textbf{দলিল} & \textbf{অবস্থা} & \textbf{প্রাসঙ্গিকতা} & \textbf{ব্লুপ্রিন্টের নকশাগত সাড়া} \\
\midrule
ব্যক্তিগত উপাত্ত সুরক্ষা আইন, ২০২৬ (Personal Data Protection Act, 2026) & ব্যক্তিগত উপাত্ত সুরক্ষা অধ্যাদেশ, ২০২৫-এর স্থলে এপ্রিল ২০২৬-এ প্রণীত বলে প্রতিবেদিত (২০২৬ সালের ৬৩ নং আইন) [7, 38]। অধ্যাদেশে প্রধান ডেটা কর্মকর্তা, অভিযোগ ও অপরাধ-সংক্রান্ত ধারাগুলো তাৎক্ষণিকভাবে কার্যকর হয়নি, এবং ১৮ মাসের সম্মতি-সময়সীমা প্রতিবেদিত হয়েছে [39, 90]। আইনের প্রতিটি বিধান ও বিধিমালার কার্যকারিতা: নিশ্চিত করতে হবে & সম্মতি; স্বয়ংক্রিয় সিদ্ধান্ত সীমিত করার অধিকারসহ ডেটা-বিষয়ীর অধিকার; সংবেদনশীল ও শিশুদের ডেটার বাড়তি সুরক্ষা; প্রধান ডেটা কর্মকর্তা; ২০২৬ সালের সংশোধনী অধ্যাদেশে পরিবর্তিত ডেটা দেশে রাখা-সংক্রান্ত বিধান [69], প্রতিবেদন অনুযায়ী গুরুত্বপূর্ণ খাতে সীমিত [70] & প্রথমে P0; যেকোনো P2 কাজের আগে DPIA; P3 বর্জিত; ব্যক্তিসংক্রান্ত কোনো স্বয়ংক্রিয় সিদ্ধান্ত নয়; বাংলাদেশে হোস্টিং \\
জাতীয় উপাত্ত ব্যবস্থাপনা আইন, ২০২৬ (National Data Management Act, 2026) & জাতীয় উপাত্ত ব্যবস্থাপনা অধ্যাদেশ, ২০২৫-এর (কিছু ইংরেজি প্রতিবেদনে National Data Governance Ordinance নামে উল্লিখিত) ধারাবাহিকতায় ১০ এপ্রিল ২০২৬-এ বিল পাস বলে প্রতিবেদিত [6, 8, 71]। কার্যকারিতা ও বিধিমালা: নিশ্চিত করতে হবে & জাতীয় উপাত্ত ব্যবস্থাপনা কর্তৃপক্ষ; দায়িত্বশীল ডেটা বিনিময় (NRDEX); জাতীয় সোর্স কোড ভান্ডার (২০২৫ সালের অধ্যাদেশের বর্ণনা অনুযায়ী [6]; আইনে বহাল আছে কি না নিশ্চিত করতে হবে) & সম্ভব হলে ডেটা প্রবেশাধিকারে NRDEX চুক্তি ব্যবহার; উন্মুক্তভাবে কোড প্রকাশ; কর্তৃপক্ষ ও সংস্থার ডেটা কর্মকর্তাদের সঙ্গে সমন্বয় \\
সাইবার নিরাপত্তা আইন, ২০২৬ ও সংশোধনী & ১০ এপ্রিল ২০২৬-এ বিল পাস বলে প্রতিবেদিত [8]; সাইবার নিরাপত্তা (সংশোধন) বিল, ২০২৬ ৩০ জুন ২০২৬-এ পাস; এর বিষয়বস্তু নিয়ে প্রতিবেদন ভিন্ন (ক্ষতিকর এআই-সৃষ্ট কনটেন্ট অপসারণ [9]; অনলাইন জুয়া-সংক্রান্ত ধারা বাদ [40]); জুলাই ২০২৬ থেকে আরেকটি খসড়া সংশোধনী মন্ত্রিসভা কমিটির পর্যালোচনাধীন [41]। বিষয়বস্তু ও কার্যকারিতা: নিশ্চিত করতে হবে & নিরাপত্তা দায়িত্ব; গুরুত্বপূর্ণ তথ্য অবকাঠামো; ক্ষতিকর কনটেন্ট অপসারণ ও বন্ধ করা, সম্ভবত এআই-সৃষ্ট কনটেন্টসহ (নিশ্চিত করতে হবে) & NIST CSF 2.0 নিয়ন্ত্রণ; ঘটনা প্রতিবেদন; তথ্য-অখণ্ডতা গবেষণা আইন ও অধিকারের সীমার মধ্যে \\
তথ্য অধিকার আইন, ২০০৯ & বলবৎ [65] & জনগুরুত্বপূর্ণ তথ্যের স্বতঃপ্রণোদিত প্রকাশ & P0 কর্পাস ও জনসমক্ষে স্বচ্ছতা নিবন্ধনকে সহায়তা করে \\
জাতীয় এআই নীতি & খসড়া ২০২৬-২০৩০ জানুয়ারি ২০২৬-এর শেষে জন-পরামর্শের জন্য প্রকাশিত [10, 11]; সংশোধিত খসড়া (সংস্করণ ২.০, ৯ ফেব্রুয়ারি ২০২৬) সরকারি পোর্টালে প্রকাশিত [12]; একটি যৌথ দলের মাধ্যমে ২০২৬-এর চতুর্থ ত্রৈমাসিকে প্রণয়ন শুরুর ঘোষণা [14]। গৃহীত হয়নি & ঝুঁকির শ্রেণি; অ্যালগরিদমিক প্রভাব মূল্যায়ন; সংসদের আইনের মাধ্যমে গঠিতব্য একটি স্বাধীন তত্ত্বাবধান কমিটি; বাংলা LLM; জাতীয় কম্পিউটিং; ২০৩০ পর্যন্ত ২০০ থেকে ২৫০ কোটি টাকার এআই উদ্ভাবন তহবিল (প্রস্তাবিত) [10] & খসড়ার শ্রেণির সঙ্গে স্তর সামঞ্জস্যপূর্ণ; ব্লুপ্রিন্ট কারিগরি উপকরণ হিসেবে প্রস্তাবিত; জাতীয় LLM পরিকল্পনার সঙ্গে BPFM সমন্বিত \\
দীর্ঘমেয়াদি জাতীয় পরিকল্পনা & প্রেক্ষিত পরিকল্পনা ২০৪১ ও আইসিটি মাস্টারপ্ল্যান ২০৪১ [15, 16] & দীর্ঘমেয়াদি ডিজিটাল রূপান্তরের লক্ষ্য & ২০৪১ দিগন্তের সঙ্গে সামঞ্জস্যপূর্ণ ও পরিকল্পনা সংশোধনের সঙ্গে খাপ খাওয়ানো যায় এমন রোডম্যাপ \\
এআই হাব ও এআই দক্ষতা প্রকল্প (কোইকা, আইসিটি বিভাগ, বাংলাদেশ হাই-টেক পার্ক কর্তৃপক্ষ) & ১৯২.৬৬ কোটি টাকার কারিগরি সহায়তা প্রকল্প প্রস্তাব জুন ২০২৬-এ সরকারি অনুমোদনের অপেক্ষায় ছিল [53]; ২০২৬ থেকে ২০২৯ মেয়াদের জন্য ১৩ মিলিয়ন মার্কিন ডলারের অনুদান-চুক্তি (রেকর্ড অব ডিসকাশন) আগস্ট ২০২৬-এ স্বাক্ষরিত বলে প্রতিবেদিত [55]। চালু হওয়ার কোনো প্রতিবেদন নেই & দক্ষতা, স্টার্টআপ, এআই গ্রহণ & পরিপূরক কর্মবিভাজন; যৌথ বেঞ্চমার্ক ও উন্মুক্ত ডেটা \\
\bottomrule
\end{tabularx}
\par\vspace{5pt}{\footnotesize\itshape অবস্থা অক্টোবর ২০২৬-এ প্রাপ্ত সংবাদমাধ্যমের প্রতিবেদন ও মাধ্যমিক আইনি ভাষ্য থেকে নেওয়া। এই সংস্করণের জন্য গেজেটভুক্ত পাঠ ও তার কার্যকারিতা-বিধান পর্যালোচনা করা হয়নি (সংযুক্তি ত), এবং বাংলাদেশের যোগ্য আইনজীবীর মাধ্যমে নিশ্চিত করতে হবে। এই সারণি আইনি পরামর্শ নয়।}
\par\vspace{3pt}{\scriptsize\color{gray} উৎস: BAPIL টেকনিক্যাল ব্লুপ্রিন্ট, সংস্করণ ১.০ (২০২৬)। CC BY 4.0।}
\end{minipage}
\end{document}
